\documentclass[10pt,a4paper]{amsart}
\usepackage[margin=19mm]{geometry}
\usepackage{amsmath,amssymb,mathtools}
\usepackage[scr=rsfs,cal=euler]{mathalfa}
\usepackage{xfrac}
\usepackage[unicode,hidelinks,bookmarks=false]{hyperref}
\newcommand{\ie}{i.e.\ }
\newcommand{\cf}{cf.\ }
\newcommand{\resp}{respectively\ }
\newcommand{\Subsection}{\subsection}
\providecommand{\nobreakdash}{\nobreak\hbox{-}}
\setlength{\emergencystretch}{2em}
\multlinegap0pt
\usepackage{bm}
\usepackage{bbm}
\usepackage{dsfont}
\usepackage{xspace}
\usepackage{cancel}
\usepackage{mathtools}
\usepackage{amsthm}
\makeatletter
\@ifundefined{theorem}{\newtheorem{theorem}{Theorem}}{}
\@ifundefined{lemma}{\newtheorem{lemma}[theorem]{Lemma}}{}
\@ifundefined{proposition}{\newtheorem{proposition}[theorem]{Proposition}}{}
\makeatother
\title[Existence of global solutions to the Fokas-Lenells equation ]{Existence of global solutions to the Fokas-Lenells equation with arbitrary spectral singularities}
\author{Yuan Li, Qiaoyuan Cheng, Engui Fan}
\date{Source: arXiv:2512.21536v1 (CC BY 4.0)}
\begin{document}
\maketitle

\noindent\emph{Source and licence.} Existence of global solutions to the Fokas-Lenells equation with arbitrary spectral singularities. Version arXiv:2512.21536v1 (CC BY 4.0), distributed under the Creative Commons Attribution 4.0 International licence, as recorded in the arXiv licence metadata for this exact version and shown on the abstract page https://arxiv.org/abs/2512.21536v1. Full rights notice: licenses/arxiv-2512.21536-CC-BY-4.0.txt.\par\smallskip

\noindent\textit{Editorial scope.} Counted unit: the Proposition whose source label is \texttt{refa} in the article (source0.tex lines 1213--1217, source label \texttt{refa}), delivered together with the article's own complete original proof of it (source0.tex lines 1218--1307, one \texttt{proof} environment of 90 lines). No mathematics is written, repaired or re-derived: statement and proof are the article's own text line for line. Required context retained uncounted: Ps (lines 432--434); timep (lines 919--921); rrh (lines 1130--1134); rrp (lines 1130--1134); rru (lines 1130--1134); po1 (lines 1191--1193); po2 (lines 1195--1197); imlem (lines 1199--1210); e1 (lines 1201--1209); e7 (lines 1201--1209). Every definition, display and label of those lines is retained unchanged. Omitted from this package and not redistributed: 1--431, 435--918, 922--1129, 1135--1190, 1194--1194, 1198--1198, 1210--1212, 1308--1449. Nothing inside a retained excerpt is omitted or altered: every retained body is delivered exactly as the article's lines are written, so no omission receipt of any kind accompanies this package. Citations: the retained excerpts carry no citation command at all, so no bibliography entry of the article is transcribed and no reference is redistributed. Reference adaptation: none was needed and none was made; every reference inside the delivered unit resolves inside it, so no unresolved reference or citation is produced. Printed slips kept literal: no printed word, symbol, hypothesis or number of the counted statement or of its proof was altered. Layout and class: the article's own class is \texttt{elsarticle} with packages including graphicx, pifont, latexsym, ifthen, amsthm, rotating, calc, textcase, booktabs, amsfonts, amssymb, amsbsy, amsmath, mathrsfs; the wrapper is the lane's pinned \texttt{amsart} editorial wrapper at 10pt on A4, which restates the theorem-like environments the retained text needs and adds the article's own macro definitions that the retained text needs (none), with extra wrapper packages (bm, bbm, dsfont, xspace, cancel, mathtools). The delivered numbering is the unit's own. Assets: no retained span contains a figure environment, a graphics-inclusion command, a PDF-page inclusion, a table or a TikZ picture; no image file is copied into this package, the package declares no asset dependency and no image file is redistributed with it. Licence: version arXiv:2512.21536v1 of this article is distributed under the Creative Commons Attribution 4.0 International licence, as recorded in the arXiv licence metadata for this exact version and shown on the abstract page https://arxiv.org/abs/2512.21536v1. A full rights notice is delivered at licenses/arxiv-2512.21536-CC-BY-4.0.txt. Characters: every retained excerpt is pure ASCII, so the delivered engine, XeLaTeX with its default Latin Modern text font, reports no missing character.
\begin{equation}\label{Ps}
\Psi_{x}=\frac{i}{2}z\left[\sigma_{3},\Psi\right]+\widetilde{U}\Psi,
\end{equation}

\begin{proposition}\label{timep}
For initial data $u_{0}\in  H^{3}(\mathbb{R})\cap H^{2,1}(\mathbb{R})$, we have $G^{(2)}_{\pm}(t,\cdot)-I \in H^{1}(\partial\widetilde{\Omega}^{\pm})\cap L^{2,1}(\partial\widetilde{\Omega}^{\pm})$ and $z^{-1}(G^{(2)}_{\pm}(t,\cdot)-I )\in H^{1}(\partial\widetilde{\Omega}^{\pm})$,  $\forall\, t\in [-T,T]$ with $T>0$. Moreover, the map $u_{0}\rightarrow G^{(2)}_{\pm}-I$ is Lipschitz continuous from  $H^{3}(\mathbb{R})\cap H^{2,1}(\mathbb{R})$ to $C([-T,T], H^{1}(\partial\widetilde{\Omega}^{\pm})\cap L^{2,1}(\partial\widetilde{\Omega}^{\pm}))$ for every $T>0$.
\end{proposition}

\begin{align}
&p(x)= \left(-\frac{1}{2\pi i} \int_{\widetilde{\Lambda}}\omega^{(2)}(x,z)\left(G^{(2)}_{x+}(z)-G^{(2)}_{x-}(z)\right)dz\right)_{12},\label{rrp}\\
&h(x)= \left(-\frac{1}{2\pi i} \int_{\widetilde{\Lambda}}\omega^{(2)}(x,z)\left(G^{(2)}_{x+}(z)-G^{(2)}_{x-}(z)\right)dz\right)_{21},\label{rrh}\\
&u(x)= -\frac{e^{-i\int_{x}^{\infty}|p(y)|^{2}dy}}{2\pi} \left(\int_{\widetilde{\Lambda}}\omega^{(2)}(x,z)\left(z^{-1}G^{(2)}_{x+}(z)-z^{-1}G^{(2)}_{x-}(z)\right)dz\right)_{12},\label{rru}
\end{align}

\begin{equation}\label{po1}
G^{(2)}_{x+}(z)-I=0,\quad z\in \widetilde{\Lambda}\setminus\widetilde{\Lambda}_{+}
\end{equation}

\begin{equation}\label{po2}
G^{(2)}_{x-}(z)-I=0,\quad z\in \widetilde{\Lambda}\setminus\widetilde{\Lambda}_{-}.
\end{equation}

\begin{lemma}\label{imlem}
Let $G^{(2)}_{\pm}-I\in H^{1}(\partial\widetilde{\Omega}^{\pm})$. Then, for all $x\geq0$, we have
\begin{align}
&\|\mathcal{C}_{\mathbb{R}}^{-}(G^{(2)}_{x+}-I)\|_{L^{2}(\mathbb{R})}\leq\frac{1}{\sqrt{1+x^{2}}}\|G^{(2)}_{+}-I\|_{H^{1}},\label{e1}\\
&\|\mathcal{C}_{\mathbb{R}}^{+}(G^{(2)}_{x-}-I)\|_{L^{2}(\mathbb{R})}\leq\frac{1}{\sqrt{1+x^{2}}}\|G^{(2)}_{-}-I\|_{H^{1}},\label{e2}\\
&\left\|\mathcal{C}_{\widetilde{\Lambda}\rightarrow\widetilde{\Lambda}_{-}}^{-}(G^{(2)}_{x+}-I)\right\|_{L^{2}(\widetilde{\Lambda}_{-})}\leq\frac{c}{\sqrt{1+x^{2}}}\|G^{(2)}_{+}-I\|_{H^{1}},\label{e3}\\
&\left\|\mathcal{C}_{\widetilde{\Lambda}\rightarrow\widetilde{\Lambda}_{+}}^{+}(G^{(2)}_{x-}-I)\right\|_{L^{2}(\widetilde{\Lambda}_{+})}\leq\frac{c}{\sqrt{1+x^{2}}}\|G^{(2)}_{-}-I\|_{H^{1}},\label{e4}\\
&\|G^{(2)}_{x+}-I\|_{L^{2}(\widetilde{\Lambda}_{2})}\leq\frac{c}{\sqrt{1+x^{2}}}\|G^{(2)}_{+}-I\|_{H^{1}},\label{e5}\\
&\|G^{(2)}_{x-}-I\|_{L^{2}(\widetilde{\Lambda}_{4})}\leq\frac{c}{\sqrt{1+x^{2}}}\|G^{(2)}_{-}-I\|_{H^{1}},\label{e6}\\
&\|(\mathcal{C}_{G^{(2)}_{x}})^{2}I\|_{L^{2}(\widetilde{\Lambda})}\leq\frac{c}{\sqrt{1+x^{2}}}\|G^{(2)}_{-}-I\|_{H^{1}}\|G^{(2)}_{+}-I\|_{H^{1}}.\label{e7}
\end{align}
\end{lemma}

\begin{proposition}\label{refa}
Suppose that $G^{(2)}_{\pm}-I\in H^{1}(\partial\widetilde{\Omega}^{\pm})\cap L^{2,1}(\partial\widetilde{\Omega}^{\pm})$
 and $z^{-1}(G^{(2)}_{\pm}-I )\in H^{1}(\partial\widetilde{\Omega}^{\pm})$.
Then the reconstructed potential satisfies $u\in H^{3}(\mathbb{R})\cap H^{2,1}(\mathbb{R})$.
\end{proposition}

\begin{proof}
Both \eqref{rrp} and \eqref{rrh} admit a decomposition
\begin{equation}\label{sgb}
\int_{\widetilde{\Lambda}}\omega^{(2)}\left(G^{(2)}_{x+}-G^{(2)}_{x-}\right)dz=F_{1}+F_{2}+F_{3}+F_{4},
\end{equation}
 with components defined as
\begin{equation*}
\begin{aligned}
&F_{1}=\int_{\widetilde{\Lambda}}\left(G^{(2)}_{x+}-G^{(2)}_{x-}\right)dz,\\
&F_{2}=\int_{\widetilde{\Lambda}}\left(\mathcal{C}_{G^{(2)}_{x}}I\right)\left(G^{(2)}_{x+}-G^{(2)}_{x-}\right)dz,\\
&F_{3}=\int_{\widetilde{\Lambda}}\left(\left(\mathcal{C}_{G^{(2)}_{x}}\right)^{2}I\right)\left(G^{(2)}_{x+}-G^{(2)}_{x-}\right)dz,\\
&F_{4}=\int_{\widetilde{\Lambda}}\left(\mathcal{C}_{G^{(2)}_{x}}\left(I-\mathcal{C}_{G^{(2)}_{x}}\right)^{-1}\left(\mathcal{C}_{G^{(2)}_{x}}\right)^{2}I\right)\left(G^{(2)}_{x+}-G^{(2)}_{x-}\right)dz.
\end{aligned}
\end{equation*}
By leveraging the properties of the Fourier and Laplace transforms, we conclude that $F_{1}\in L^{2,1}$. Since $F_{2}$ is a diagonal matrix, it does not contribute to the reconstruction of $p$ and $h$.

The estimation of $F_{3}$ and $F_{4}$ requires careful application of both the definition of the Cauchy integral operators $\mathcal{C}_{\widetilde{\Lambda}\rightarrow\widetilde{\Lambda}'}^{\pm}$ and their boundedness between the function spaces $L^{2}(\widetilde{\Lambda})$ and $L^{2}(\widetilde{\Lambda}')$. From \eqref{po1}--\eqref{po2}, we obtain
\begin{equation}
\begin{aligned}\nonumber
\left|F_{3}\right|\leq&\left|\left(\int_{\mathbb{R}}+\int_{\widetilde{\Lambda}_{2}}\right)\left(\mathcal{C}_{\widetilde{\Lambda}}^{+}\left(\mathcal{C}_{\widetilde{\Lambda}}^{-}\left(G^{(2)}_{x+}-I\right)\right)\left(I-G^{(2)}_{x-}\right)\right)\left(G^{(2)}_{x+}-I\right)dz\right|\\
&+\left|\left(\int_{\mathbb{R}}+\int_{\widetilde{\Lambda}_{4}}\right)\left(\mathcal{C}_{\widetilde{\Lambda}}^{-}\left(\mathcal{C}_{\widetilde{\Lambda}}^{+}\left(I-G^{(2)}_{x-}\right)\right)\left(G^{(2)}_{x+}-I\right)\right)\left(I-G^{(2)}_{x-}\right)dz\right|\\
\leq&\left|\int_{\mathbb{R}}\left(\mathcal{C}_{\widetilde{\Lambda}_{-}\rightarrow\mathbb{R}}^{+}\left(\mathcal{C}_{\widetilde{\Lambda}\rightarrow\widetilde{\Lambda}_{-}}^{-}\left(G^{(2)}_{x+}-I\right)\right)\left(I-G^{(2)}_{x-}\right)\right)\mathcal{C}_{\mathbb{R}}^{-}\left(G^{(2)}_{x+}-I\right)dz\right|\\
&+\left|\int_{\widetilde{\Lambda}_{2}}\left(\mathcal{C}_{\widetilde{\Lambda}_{-}\rightarrow\widetilde{\Lambda}_{2}}^{+}\left(\mathcal{C}_{\widetilde{\Lambda}\rightarrow\widetilde{\Lambda}_{-}}^{-}\left(G^{(2)}_{x+}-I\right)\right)\left(I-G^{(2)}_{x-}\right)\right)\left(G^{(2)}_{x+}-I\right)dz\right|\\
&+\left|\int_{\mathbb{R}}\left(\mathcal{C}_{\widetilde{\Lambda}_{+}\rightarrow\mathbb{R}}^{-}\left(\mathcal{C}_{\widetilde{\Lambda}\rightarrow\widetilde{\Lambda}_{+}}^{+}\left(I-G^{(2)}_{x-}\right)\right)\left(G^{(2)}_{x+}-I\right)\right)\mathcal{C}_{\mathbb{R}}^{+}\left(I-G^{(2)}_{x-}\right)dz\right|\\
&+\left|\int_{\widetilde{\Lambda}_{4}}\left(\mathcal{C}_{\widetilde{\Lambda}_{+}\rightarrow\widetilde{\Lambda}_{4}}^{-}\left(\mathcal{C}_{\widetilde{\Lambda}\rightarrow\widetilde{\Lambda}_{+}}^{+}\left(I-G^{(2)}_{x-}\right)\right)\left(G^{(2)}_{x+}-I\right)\right)\left(I-G^{(2)}_{x-}\right)dz\right|.
\end{aligned}
\end{equation}
By applying estimates \eqref{e1}--\eqref{e7} established in Lemma \ref{imlem}, we can further derive the following bound for $\left|F_{3}\right|$
\begin{align*}
\left|F_{3}\right|\lesssim&\left\|\left(\mathcal{C}_{\widetilde{\Lambda}\rightarrow\widetilde{\Lambda}_{-}}^{-}\left(G^{(2)}_{x+}-I\right)\right)\right\|_{L^{2}(\widetilde{\Lambda}_{-})}\left(\left\|\mathcal{C}_{\mathbb{R}}^{-}\left(G^{(2)}_{x+}-I\right)\right\|_{L^{2}(\mathbb{R})}+\left\|G^{(2)}_{x+}-I\right\|_{L^{2}(\widetilde{\Lambda}_{2})}\right)\\
&+\left\|\left(\mathcal{C}_{\widetilde{\Lambda}\rightarrow\widetilde{\Lambda}_{+}}^{+}\left(I-G^{(2)}_{x-}\right)\right)\right\|_{L^{2}(\widetilde{\Lambda}_{+})}\left(\left\|\mathcal{C}_{\mathbb{R}}^{+}\left(I-G^{(2)}_{x-}\right)\right\|_{L^{2}(\mathbb{R})}+\left\|I-G^{(2)}_{x-}\right\|_{L^{2}(\widetilde{\Lambda}_{4})}\right)\\
\lesssim&\,\frac{1}{1+x^{2}}.
\end{align*}
The fourth integral component $F_{4}$ admits the representation
\begin{equation*}
F_{4}=\int_{\widetilde{\Lambda}}\left(\mathcal{C}_{G^{(2)}_{x}}f\right)\left(G^{(2)}_{x+}-G^{(2)}_{x-}\right)dz,
\end{equation*}
with $f=(I-\mathcal{C}_{G^{(2)}_{x}})^{-1}(\mathcal{C}_{G^{(2)}_{x}})^{2}I$. The boundedness property of $(I-\mathcal{C}_{G^{(2)}_{x}})^{-1}$ combined with \eqref{e7} establishes the estimate
\begin{equation*}
\|f\|_{L^{2}(\widetilde{\Lambda})}\lesssim\|(\mathcal{C}_{G^{(2)}_{x}})^{2}I\|_{L^{2}(\widetilde{\Lambda})}\lesssim\frac{1}{\sqrt{1+x^{2}}}.
\end{equation*}
It follows from \eqref{po1}--\eqref{po2} that
\begin{equation*}
\begin{aligned}
\left|F_{4}\right|\leq&\left|\left(\int_{\mathbb{R}}+\int_{\widetilde{\Lambda}_{2}}\right)\left(\mathcal{C}_{\widetilde{\Lambda}}^{+}f\left(I-G^{(2)}_{x-}\right)\right)\left(G^{(2)}_{x+}-I\right)dz\right|\\
&+\left|\left(\int_{\mathbb{R}}+\int_{\widetilde{\Lambda}_{4}}\right)\left(\mathcal{C}_{\widetilde{\Lambda}}^{-}f\left(G^{(2)}_{x+}-I\right)\right)\left(I-G^{(2)}_{x-}\right)dz\right|\\
\leq&\left|\int_{\mathbb{R}}\left(\mathcal{C}_{\widetilde{\Lambda}\rightarrow\mathbb{R}}^{+}f\left(I-G^{(2)}_{x-}\right)\right)\mathcal{C}_{\mathbb{R}}^{-}\left(G^{(2)}_{x+}-I\right)dz\right|\\
&+\left|\int_{\widetilde{\Lambda}_{2}}\left(\mathcal{C}_{\widetilde{\Lambda}\rightarrow\widetilde{\Lambda}_{2}}^{+}f\left(I-G^{(2)}_{x-}\right)\right)\left(G^{(2)}_{x+}-I\right)dz\right|\\
&+\left|\int_{\mathbb{R}}\left(\mathcal{C}_{\widetilde{\Lambda}\rightarrow\mathbb{R}}^{-}f\left(G^{(2)}_{x+}-I\right)\right)\mathcal{C}_{\mathbb{R}}^{+}\left(I-G^{(2)}_{x-}\right)dz\right|\\
&+\left|\int_{\widetilde{\Lambda}_{4}}\left(\mathcal{C}_{\widetilde{\Lambda}\rightarrow\widetilde{\Lambda}_{4}}^{-}f\left(G^{(2)}_{x+}-I\right)\right)\left(I-G^{(2)}_{x-}\right)dz\right|.
\end{aligned}
\end{equation*}
Then, an application of Lemma \ref{imlem} yields
\begin{equation*}
\begin{aligned}
\left|F_{4}\right|
\lesssim&\left\|f\left(I-G^{(2)}_{x-}\right)\right\|_{L^{2}(\widetilde{\Lambda})}\left(\left\|\mathcal{C}_{\mathbb{R}}^{-}\left(G^{(2)}_{x+}-I\right)\right\|_{L^{2}(\mathbb{R})}+\left\|G^{(2)}_{x+}-I\right\|_{L^{2}(\widetilde{\Lambda}_{2})}\right)\\
&+\left\|f\left(G^{(2)}_{x+}-I\right)\right\|_{L^{2}(\widetilde{\Lambda})}\left(\left\|\mathcal{C}_{\mathbb{R}}^{+}\left(I-G^{(2)}_{x-}\right)\right\|_{L^{2}(\mathbb{R})}+\left\|I-G^{(2)}_{x-}\right\|_{L^{2}(\widetilde{\Lambda}_{4})}\right)\\
\lesssim&\,\frac{1}{1+x^{2}}.
\end{aligned}
\end{equation*}
Base on the above analysis for $F_{i}$, $1\leq i\leq4$, we conclude that $p,h\in L^{2,1}(\mathbb{R})$. Proposition \ref{timep} implies that $z^{-1}(G^{(2)}_{\pm}-I )\in H^{1}(\partial\widetilde{\Omega}^{\pm})$. Thus, by repeating the same analysis to \eqref{rru}, we can obtain $u\in L^{2,1}(\mathbb{R})$.


Finally, we establish that $h_{x}\in L^{2}(\mathbb{R})$. Defining the transform $L(x,z)=e^{i\int_{x}^{\infty}|u_{y}(y)|^{2}dy\sigma_{3}}\Psi(x,z)$, it follows from \eqref{Ps} that
\begin{equation}\label{ahf}
L_{x}=\frac{i}{2}z\left[\sigma_{3},L\right]+\widetilde{Q}L,
\end{equation}
where
\begin{equation*}
\widetilde{Q}=e^{i\int_{x}^{\infty}|u_{y}(y)|^{2}dy\widehat{\sigma}_{3}}\begin{pmatrix}0&iu_{x}\\
-\bar{u}_{xx}-i\bar{u}_{x}|u_{x}|^{2}&0\end{pmatrix}.
\end{equation*}
Moreover, $\omega^{(2)}$ satisfies \eqref{ahf}, leading to the derivative relation
\begin{equation*}
\frac{d}{dx}\left(\omega^{(2)} \left(G^{(2)}_{x+}-G^{(2)}_{x-}\right)\right)=\frac{i}{2}z\left[\sigma_{3},\omega^{(2)} \left(G^{(2)}_{x+}-G^{(2)}_{x-}\right)\right]+\widetilde{Q}\,\omega^{(2)} \left(G^{(2)}_{x+}-G^{(2)}_{x-}\right).
\end{equation*}
By integrating both sides of this equation, we derive
\begin{equation}
\frac{d}{dx}\int_{\widetilde{\Lambda}}\omega^{(2)} \left(G^{(2)}_{x+}-G^{(2)}_{x-}\right)dz=I_{1}(x)+I_{2}(x),
\end{equation}
where
\begin{equation*}
\begin{aligned}
&I_{1}(x)=\frac{i}{2}\int_{\widetilde{\Lambda}}\left[\sigma_{3},\omega^{(2)} z\left(G^{(2)}_{x+}-G^{(2)}_{x-}\right)\right]dz,\\
&I_{2}(x)=\int_{\widetilde{\Lambda}}\widetilde{Q}\,\omega^{(2)} \left(G^{(2)}_{x+}-G^{(2)}_{x-}\right)dz.
\end{aligned}
\end{equation*}
The term $I_{1}$ admits a decomposition analogous to \eqref{sgb}, yielding $I_{1}\in L^{2}(\mathbb{R})$  under the condition that $z(G^{(2)}_{\pm}-I)\in L^{2}(\partial \widetilde{\Omega}^{\pm})$. Analysis of \eqref{sgb} establishes $\widetilde{Q}\in L^{\infty}(\mathbb{R})$, which consequently implies $I_{2}\in L^{2}(\mathbb{R})$. These results collectively demonstrate that $h_{x}\in L^{2}(\mathbb{R})$.
\end{proof}

\end{document}
